\documentclass[times,twocolumn,final]{elsarticle}

\usepackage{cag}
\usepackage{framed,multirow}
\usepackage{amssymb}
\usepackage{amsmath}
\usepackage{lineno}
\usepackage{graphicx}
\usepackage{booktabs}
\setcounter{totalnumber}{6}
\setcounter{topnumber}{4}
\setcounter{bottomnumber}{4}
\renewcommand{\topfraction}{0.92}
\renewcommand{\dbltopfraction}{0.95}
\renewcommand{\textfraction}{0.05}
\renewcommand{\floatpagefraction}{0.85}


\journal{Computers \& Graphics}

\begin{document}

\begin{frontmatter}

\title{Timestep-Conditioned Adapter Scaling for Multi-view Diffusion Texture Generation}

% Double-anonymized review: author names, affiliations, and acknowledgements
% (including funding) are provided only on the separately uploaded title page.

\begin{abstract}
Geometry-conditioned multi-view diffusion relies on inference-time adapter strength to balance geometric alignment and generated appearance. We study Timestep-Conditioned Adapter Scaling (TCAS), a training-free temporal control on the existing GeoTex-Adapter residual path. The low--high--low profile C3 is retained from a 24-object development study. In the original blinded study, C3 received 58.1\% overall-quality preference (95\% cluster-bootstrap CI $[50.1,65.8]$) across 30 objects and 24 participants. A new paired evaluation on 300 input-disjoint Fresh C objects gives C3--Fixed Low FG-PSNR $+0.52$ dB (95\% object-bootstrap CI $[0.40,0.64]$), while stronger fixed and linear controls have better means on the four image endpoints. Four equal-applied-budget controls within a six-condition addendum distinguish stage allocation from strength: C3 trails Fixed Mean and Late-high by 0.13 and 0.56 dB FG-PSNR. On ten previously unbaked, compatibility-selected objects, calibrated native-UV baking gives Late-high--Fixed Mean unseen FG-LPIPS $-0.00265$ ($[-0.00414,-0.00137]$), favorable on all ten, although substantial reference-appearance errors remain. A fifty-object native MV-Adapter comparison gives Middle-high--Fixed Mean FG-PSNR $+0.029$ dB ($[-0.019,+0.081]$). The historical preference, paired image comparisons and downstream evaluation characterize adapter strength as a measurable stage-dependent control. The supported benefits are specific to their checkpoint, comparator and endpoint; temporal scheduling itself is established prior work, and the evidence does not establish a universally preferred middle window.
\end{abstract}

\begin{keyword}
\KWD
multi-view diffusion \sep 3D texturing \sep geometry adapter \sep residual scaling \sep timestep conditioning
\end{keyword}

\end{frontmatter}

\linenumbers

\section{Introduction}
\label{sec:intro}

Multi-view texture generation uses image diffusion to synthesize target views and then projects or bakes them onto an existing surface. Geometry adapters make normals, depth and foreground information available to the denoiser, but the strength of that information also affects surface appearance. An aggregate structural score alone cannot determine a useful setting: local variation may represent detail or unwanted patterns, and image-space agreement must be considered alongside the final projected appearance.

The 01549 study investigated this inference-control problem in an MVPainter-style pipeline through fixed-scale diagnostics, a 24-object comparison of 17 control variants, and a blinded preference study. TCAS applies a requested scale over denoising progress without training additional parameters or modifying the base model or adapter weights. Its recorded C3 profile, $(1.25,2.50,1.25)$, uses a stronger middle request and conservative early/late requests. This revision retains that method and development identity and evaluates the profile against stronger controls rather than searching for another C3 configuration.

The central question is when a temporal control improves the tested outcome beyond an available fixed scale. We distinguish the requested multiplier, its applied value after native caps, and the resulting activation-dependent correction. A matched 300-object Fresh C evaluation supplies common-input paired comparisons; a same-cohort four-control budget match separates stage allocation from a larger applied-scale sum. Calibrated unseen-view baking and a native-camera MV-Adapter test extend the evidence to downstream appearance and a second conditioning interface. The original preference study remains a separate subjective evaluation of its own stimuli.

Layer- and timestep-dependent scheduling is established prior work. Scheduled Style Injection (SSI) includes layer/time schedules and geometry-conditioned ControlNet strength scheduling. TCAS's contribution here is its concrete application and evaluation on the GeoTex residual path for multi-view texturing, not a new general scheduling primitive. The original preference and C3--Fixed Low gain support specific benefits; the stronger controls and second-interface results define their scope.

The contributions are:

(1) \textbf{Executable inference-time control.} TCAS specifies a temporal scale on an existing geometry adapter, with explicit requested/applied semantics and no additional learned parameters. Historical fixed-scale and C3 development records explain the method's origin.

(2) \textbf{Matched comparison with strong controls.} The 300-object Fresh C study reports paired gains, endpoint trade-offs and object-level variation. Its six-condition addendum combines four equal-applied-budget allocation controls with two strength controls, distinguishing the stage question from the C3--Fixed Low amplitude change (Table~\ref{tab:matched-budget300}).

(3) \textbf{Preference, downstream and second-interface evidence.} The original blinded preference study supports its overall-quality finding. Calibrated eleven-unseen-view baking on ten objects supplies a named Late-high--Fixed Mean LPIPS benefit, while native MV-Adapter testing on fifty objects delineates transfer limits (Tables~\ref{tab:preference}, \ref{tab:new-bake} and \ref{tab:new-mva}). Each result retains its own cohort, checkpoint and endpoint identity.

\section{Related Work}
\label{sec:related}

\subsection{3D Texture Generation and Multi-view Diffusion}
\label{subsec:3dtex}

Single-image 3D generation is often decomposed into geometry generation and texture generation. Single-image-to-3D, multi-view diffusion, and sparse-view reconstruction methods have improved structural stability and view synthesis \cite{liu2023zero1to3,shi2023zero123pp,liu2023syncdreamer,long2023wonder3d,shi2024mvdream,wang2023imagedream,xu2024instantmesh,ke2024era3d}. Texture synthesis remains a distinct challenge for asset realism and usability \cite{perla2026survey}. Many systems generate RGB views with 2D diffusion and project or bake them onto a surface \cite{chen2023text2tex,richardson2023texture,poole2023dreamfusion,chen2023fantasia3d}. Recent work also explores exemplar-based texture synthesis and diffusion-based procedural materials \cite{guzman2026adaptive,lv2025interpretable}.

Paint3D, MVPaint, Hunyuan3D 2.0, and MVPainter use geometry-conditioned multi-view generation in 3D texturing pipelines \cite{zeng2024paint3d,cheng2024mvpaint,lai2025hunyuan3d,shao2025mvpainter}. MV-Adapter, MVDiffusion, and other synchronized multi-view methods address consistency through different adapter and denoising designs \cite{huang2024mvadapter,tang2023mvdiffusion,liu2024synchronized,xu2026wondertex}. These systems establish the graphics context for this work. The present evidence does not compare final baked-texture quality across those systems.

\subsection{Adapter-Augmented Diffusion Models and Geometric Control}
\label{subsec:adapter}

ControlNet and T2I-Adapter provide learned spatial-conditioning branches or feature injections, whose contribution can be weighted by an inference-time condition-strength scale \cite{mou2023t2iadapter,zhang2023controlnet}. They establish conditioning strength as a practical control, but their injected tensors and scaling sites do not share a numerical dose with GeoTex. The GeoTex-Adapter path in the 01549 study already exposes an inference-time residual scale and evaluates fixed requests. TCAS reuses that control point as a three-stage request over sampling progress, while recording how wrapper-group caps change the applied values and how adapter activations determine the resulting corrections. It adds no learned branch or schedule selector. IP-Adapter introduces image prompts through decoupled attention \cite{ye2023ipadapter}; LoRA and related methods use scaled low-rank updates \cite{hu2022lora,liu2024dora,zhang2023adalora}.

\subsection{Diffusion Guidance Scheduling and Timestep Control}
\label{subsec:guidance}

Diffusion denoising is often discussed in terms of coarse-to-fine evolution, and prior work varies guidance strength over sampling progress to study quality, diversity, or fidelity \cite{ho2020ddpm,song2021ddim,ho2022cfg,castillo2023adaptive}. Limited-interval classifier-free guidance is one example; it changes the combination of conditional and unconditional score predictions \cite{kynkanniemi2024limited}.

SSI is the closest scheduling precedent. It varies StyleID's query-mixing parameter across decoder layers and denoising time, tests schedule direction and shape, and applies scheduling to a separate geometry-conditioned ControlNet branch; its reported ControlNet scheduling gains over fixed conditioning are modest relative to its style-injection results \cite{kulkarni2026ssi}. This prior work already covers the general idea of layer/time scheduling and geometry-conditioning schedules. TCAS acts on the residual correction produced by GeoTex-Adapter from normal, depth, and foreground-mask conditions, with group-specific wrapper caps; it does not modulate StyleID attention mixing or add a ControlNet branch. The task, interface, and measured endpoints differ, so SSI does not imply the Fresh C paired estimates, the requested/applied/residual accounting, or the observed image- and color-agreement failures. These task-specific measurements are the empirical contribution; the scheduling concept itself is not claimed as new.

\subsection{Multi-view Adapter Interfaces}
\label{subsec:mvinterfaces}

MV-Adapter uses a condition guider for camera or geometry maps together with decoupled spatial, multi-view, and optional image cross-attention layers \cite{huang2024mvadapter}. MVDiffusion uses correspondence-aware attention modules to connect pixels across views \cite{tang2023mvdiffusion}. These interfaces differ from the GeoTex-Adapter residual branch studied here. We therefore report each auxiliary intervention within its own tested implementation and cohort, without pooling metrics or interpreting it as an architecture-matched replication of TCAS. The main cross-interface results are summarized in Table~\ref{tab:cross-interface}.

\section{Method}
\label{sec:method}

\subsection{Task Definition and Base Pipeline}
\label{subsec:task}

This paper builds on the MVPainter-style geometry-controlled multi-view texture generation framework \cite{shao2025mvpainter} and studies how the inference-time strength of geometric condition injection affects the generated results. Figure~\ref{fig:overview} shows the TCAS control gate on the GeoTex-Adapter residual path: the requested scale is limited by the wrapper's group-specific cap before the residual is added to frozen UNet features. Given a single reference image and the corresponding 3D geometry, the base pipeline generates RGB texture images from six target views and can pass them to downstream projection or texture baking. Fresh C evaluates target-view images; the separate ten-object pilot evaluates calibrated unlit baked meshes. We do not modify the base UNet, VAE, or pretrained adapter weights.

\begin{figure*}[t]
\centering
\includegraphics[width=\textwidth]{figures/fig1_tcas_control.pdf}
\caption{TCAS on the existing GeoTex residual path. Reference conditioning and five-channel geometry conditioning enter separate paths. The requested scale is capped per layer group before injection into the frozen UNet. The diagram abstracts multiple ResNet sites; steps are zero-based, and Fresh C shallow requests are capped at 0.8. This is an implementation schematic, not an experimental result.}
\label{fig:overview}
\end{figure*}

During multi-view diffusion, the 3D geometry is rendered into normal maps, depth maps, and foreground masks, denoted as:
\begin{equation}
\mathbf{G} = \{G_\text{normal}, G_\text{depth}, G_\text{mask}\}.
\end{equation}

The geometric adapter maps these conditions to a residual correction term and injects it into the intermediate hidden states of the UNet. For inference step $j$ with normalized execution progress $p_j$, the adapter injection can be written as:
\begin{equation}
\mathbf{h}'_{j,\ell} = \mathbf{h}_{j,\ell} + s_{\mathrm{app}}(p_j,\ell)\, A^\ell_\phi(\mathbf{h}_{j,\ell}, \mathbf{G}),
\label{eq:injection}
\end{equation}
Here $\ell$ indexes an injected layer group. The \emph{requested scale} $s_{\mathrm{req}}(p_j,\ell)$ is the value assigned by the profile; the \emph{applied scale} is the value used by the wrapper,
\begin{equation}
s_{\mathrm{app}}(p_j,\ell)=\min\!\bigl(s_{\mathrm{req}}(p_j,\ell),c_\ell\bigr),
\end{equation}
when the native cap $c_\ell$ is active, and otherwise equals the request. The resulting correction is $\mathbf{r}_{j,\ell}=s_{\mathrm{app}}(p_j,\ell)A^\ell_\phi(\mathbf{h}_{j,\ell},\mathbf{G})$; its \emph{post-scale residual magnitude} is $m_{j,\ell}=\|\mathbf{r}_{j,\ell}\|_2$. Thus neither a requested nor an applied scale label alone determines the correction magnitude.

In the audited Fresh C wrapper, the three cap groups map the GeoTex residual wrappers on the UNet ResNet blocks as follows: deep comprises the UNet middle block and $\texttt{up\_0}$, middle is $\texttt{up\_1}$, and shallow is $\texttt{up\_2}$. These names identify the implementation's grouped injection points; they do not assert universal semantic roles for denoising stages. The corresponding Fresh C caps are 3.0, 3.5, and 0.8. C3's shallow request of 1.25 is therefore applied as 0.8, while its deep and middle requests pass unchanged. GFL and GFH also clip at the shallow cap. These equations and caps describe the explicitly requested-scale path. The unset-scale GeoTex wrapper bypasses this cap operation and is a distinct path; its generated outputs are not available here (Supplementary Section S20). These cap values describe the audited Fresh C wrapper; row-level cap semantics for the historical 24-object probe are not fully recoverable. TCAS scales a geometry residual added to feature activations; it does not reweight conditional and unconditional score predictions as classifier-free guidance does.

\subsection{Image-Space Metrics and Texture-Variation Diagnostics}
\label{subsec:metrics}

We use separate terms for separate outcomes. \emph{Image-space agreement} refers to target-view PSNR, LPIPS, and SSIM measurements on the named image regions; it is not a measure of final 3D material quality. \emph{Reference-color agreement} refers to color differences such as CIEDE2000 computed against the corresponding reference. \emph{Texture variation} refers to statistics such as RGB standard deviation, Laplacian variance, gradient magnitude, and high-frequency energy; these may rise with detail or artifacts and are not fidelity measures. \emph{Subjective preference} refers only to the submitted forced-choice study. \emph{Baked 3D appearance} refers to rendered mesh outputs after texture projection or baking, while \emph{PBR/material fidelity} additionally concerns appearance under the material and illumination model. These endpoints are not interchangeable. FG-SSIM, Edge-SSIM, and Crop-SSIM quantify structural similarity on the indicated image regions; PSNR measures pixel error; LPIPS compares learned image features \cite{wang2004ssim,zhang2018lpips}.

Laplacian variance, RGB standard deviation, and gradient magnitude can describe local high-frequency or color variation. For example, a foreground Laplacian statistic may be written as
\begin{equation}
M_{\mathrm{var}}=\operatorname{Var}\!\left(\operatorname{Lap}(I_{\mathrm{out}}\odot G_{\mathrm{mask}})\right),
\end{equation}
where $G_{\mathrm{mask}}$ is the foreground mask. A larger value is not necessarily better: detail, noise, repeated patterns, or artifacts can all increase local variation. We therefore do not interpret these proxies as texture fidelity or use them alone to claim that a schedule better preserves reference appearance.

\subsection{TCAS Control Semantics and Comparison Design}
\label{subsec:tcas}

TCAS replaces a globally fixed requested adapter scale with a piecewise profile over normalized sampling progress. The early/middle/late partition is an operational description, not a universal claim about the function of every denoising stage. C3's low--high--low profile is the specific development choice evaluated in this revision; it is not presented as an independently derived optimum.

To avoid ambiguity caused by different raw timestep directions in different schedulers, $p_j=j/(T-1)$ denotes execution progress for inference step $j\in\{0,\ldots,T-1\}$. Thus $p=0$ is the first update from the initial noisy latent and $p=1$ is the last denoising update before decoding. For $T=50$, TCAS divides execution into steps 0--16, 17--32, and 33--49 (17/16/17 steps):
\begin{equation}
s_{\mathrm{req}}(p) = \begin{cases} s_e, & 0 \leq p < 1/3 \\ s_m, & 1/3 \leq p < 2/3 \\ s_l, & 2/3 \leq p \leq 1 \end{cases}.
\label{eq:tcas}
\end{equation}

We call this variation across the 50 inference steps the \emph{temporal allocation}. A \emph{layer allocation} varies the requested scale across the three wrapper groups at a fixed step. For a schedule with applied values $s_{\mathrm{app}}(j,\ell)$, its per-group applied-scale sum is $B_\ell=\sum_j s_{\mathrm{app}}(j,\ell)$; the unweighted sum $B_{\mathrm{nom}}=\sum_\ell B_\ell$ is a \emph{nominal applied-scale budget}. It does not weight scales by adapter activations. We log post-scale correction summaries, including wrapper-level $L_2$ norm and mean absolute entry. The three wrapper groups operate on feature maps with different spatial dimensions, and the traces do not retain channel counts; raw $L_2$ values are therefore not used for cross-layer dose comparisons. The per-entry mean absolute correction reported in Supplementary Section~S4 is a dimension-normalized descriptive summary, not an equal-dose intervention. Correction magnitudes also depend on the current hidden state and earlier denoising updates, so equal nominal budgets do not imply equal realized corrections.

The selected C3 configuration used in this paper is:
\begin{equation}
(s_e, s_m, s_l) = (1.25, 2.50, 1.25).
\end{equation}

The low--high--low profile is a concrete inference-time intervention motivated by the submitted development analysis. Its stage interpretation is a task-specific hypothesis: the tested early, middle, and late allocations may produce different responses. Fresh C evaluates the recorded profile against matched controls but does not prove a universal function for any stage. TCAS adds no trained parameters and does not modify base-model or adapter weights; it replaces the requested scale during inference.

\paragraph{What C3 changes and what its comparisons identify}
C3 requests 1.25 during steps 0--16, 2.50 during steps 17--32, and 1.25 during steps 33--49, applying that temporal profile to each of the three layer groups before native caps. In Fresh C, the applied deep and middle profiles retain those values; the shallow profile is clipped to 0.8 throughout. GFL requests 1.25 at every step and group, so its deep and middle values are 1.25 throughout and its shallow value is also clipped to 0.8. Their per-group applied-scale sums are therefore 82.5 versus 62.5 in deep and middle, and 40 for both in shallow. The C3--GFL contrast changes the middle-third deep/middle scale magnitude as well as its temporal allocation; it is not a timing-only comparison.

To identify a temporal-placement contrast, controls must hold the per-layer applied values and nominal per-group sums fixed while moving the same values between time windows. To identify a scale-strength contrast, controls must hold the layer and temporal allocations fixed while changing amplitude. Both contrasts also require reporting post-scale correction summaries: matching requested values or applied-scale sums does not match the evolving residual corrections. Raw $L_2$ values are dimension-sensitive across wrapper groups and cannot establish cross-layer dose matching. The current 300-object allocation comparison exactly matches the applied sums in each layer group. This controls the scalar budget for the total allocation effect, although it does not identify a residual-dose-independent mechanism. Post-hoc adjustment for state-dependent corrections cannot by itself recover that mechanism. The older strict-276 analysis is reported separately.

\paragraph{Development origin and scope of C3}
The submitted 01549 paper records C3 $=(1.25,2.50,1.25)$ as the selected profile from a 24-object development probe containing 17 uniform, layer-wise, and time-conditioned variants. We retain that development record because it explains the origin of the profile studied here. The complete row-level outputs and selection-time provenance are not available in the current evidence authority, so the historical ranking and labels cannot be independently checked. We therefore make no formal optimality or uniqueness claim based on this probe. Fresh C evaluates the recorded profile without further schedule search.

\subsection{Geometry Adapter and Inference Interface}
\label{subsec:adapterimpl}

GeoTex-Adapter is treated as the fixed geometry-residual backbone, not as a new contribution of this paper. It receives a five-channel input comprising a three-channel normal map, a one-channel depth map, and a one-channel foreground mask. A multi-scale encoder produces residual features that are injected into the frozen UNet. TCAS acts on the requested multiplier along denoising progress; it does not retrain the adapter, the UNet, or the VAE and introduces no learned parameters.

The first-round 01549 revision also reported a learned Full Adaptive Correction (FAC) extension. We retain its controlled negative re-examination as a previously reported auxiliary result in Section~\ref{subsec:facreported} and Supplementary Table~S16; it is not a new 1008 training run. TCAS remains the paper's sole training-free method.

\section{Experiments}
\label{sec:experiments}

\subsection{Experimental Setup}
\label{subsec:setup}

\subsubsection{Experimental protocol and comparisons}

The original 01549 submission reports fixed-scale experiments, a 24-object schedule-development probe, and a 300-object pooled evaluation. The source figures and tables are authentic submission artifacts, but the current evidence archive does not contain a complete object-level manifest and run record for every historical result. We therefore preserve the fixed-scale and C3-selection questions as historical development evidence and do not present their unrecovered numbers as newly verified confirmatory results.

Fresh C is a held-out, matched image-space study on 300 frozen object identifiers, independent of the original 24-object schedule-selection probe and the original 300-object evaluation cohort. The primary four-condition run (No Adapter, GFL, GFH, and C3) and the same-cohort Generic Linear/LLH addendum use the same frozen adapter checkpoint (identity recorded in Supplementary Section S8), runner source (identity recorded in Supplementary Section S8), 300-object cohort, reference/geometry inputs, target views, metric implementation, and per-object seed policy $42+\texttt{object\_idx}$. The verified gates cover 1,200/1,200 primary rows and 600/600 addendum rows. The registered primary endpoint is paired C3--GFL FG-PSNR. Fresh C's 300 UIDs have zero overlap with the original evaluation cohort (including its probe), the original 1,118-object training list, and the clean-v2 evaluation cohort. The Fresh C comparison uses one frozen checkpoint and common execution protocol; hash identity with the historical Table~4 checkpoint is not used as a validity condition. The historical distinction remains: 01549 names a v1 checkpoint for its main tables and a separate v2 strong-residual, per-layer-capped adapter for Section~4.7/FAC; Fresh C is a controlled retraining, and its estimate is not presented as a numerical reproduction of the old $+0.96$~dB contrast. The separate post-inspection strict-276 analyses use the clean-v2 object list and record the same checkpoint SHA as Fresh C; their cohort and run provenance differ from Fresh C, and they are reported only as bounded supporting diagnostics.

\begin{table*}[t]\centering\scriptsize
\caption{Evidence cohorts used for current conclusions. Historical development, preference and FAC records retain their original identities and are not pooled with these rows.}
\label{tab:cohort-map}
\begin{tabular}{llll}
\toprule Cohort & Objects / controls & Evaluation domain & Role \\
\midrule
Fresh C primary + addendum & 300 / 6 & GeoTex target images & UID-disjoint matched checkpoint comparison \\
Equal-budget addendum & Same 300 / 6 & GeoTex target images & Post-inspection allocation comparison \\
New native MV-Adapter & 50 / 7 & Native-camera albedo targets & Internally frozen second-interface test \\
New bake pilot & 10 / 7 + GT & Eleven unseen views per object & Input-compatible pilot; image outcomes previously known \\
Corrected earlier bake & Separate 5 & Unlit unseen views & Descriptive; supersedes erroneous estimates \\
\bottomrule\end{tabular}
\end{table*}

\subsubsection{Evaluation sets and metrics}

The submitted paper describes a 50-object, 13-scale sweep; a purposive 26-object texture audit; a 24-object, 17-variant development probe; and a pooled 300-object evaluation. We retain the 50-object sweep and the 24-object probe as historical diagnostics because they motivate the scale-allocation question. Their object-level raw outputs, full selection provenance, or metric implementation records are incomplete in the current archive. The 26-object audit's object IDs, selection trace, and proxy definitions also cannot be reconstructed; its prevalence and texture-loss claims are withdrawn.

Fresh C adds a revision-era 300-object image-space comparison with foreground PSNR, LPIPS, SSIM, and Edge-SSIM. These measure agreement with target-view renderings on their named endpoints; they do not measure texture fidelity, material appearance, seams, or human preference. The separately reported 01549 preference study is a subjective choice task, not a texture-fidelity measure. Laplacian variance, RGB standard deviation, and gradient magnitude are treated only as variation diagnostics, since larger values can also reflect artifacts. Object is the inferential unit; views are averaged within object and are not independent samples.

\subsubsection{Implementation and dataset details}

The authenticated 01549 source names \texttt{geotex\_refattn\_v1/geotex\_step\_0002000.pt} for its main tables but provides no checkpoint hash. Fresh C uses a controlled-retraining GeoTex-Adapter checkpoint, trained from the same 1,118-object UID list (identity recorded in Supplementary Section S8). The primary run and same-cohort addendum share this frozen checkpoint, runner source, seed policy, inputs, and endpoint implementation. Fresh C uses the MVPainter-style joint six-view latent UNet and VAE, EulerDiscreteScheduler for 50 steps, six fixed target views at 256$\times$256, and per-object seed $42+\texttt{object\_idx}$, shared across all compared conditions for that object. The 300-object manifest, per-condition row counts, and input/output hashes are indexed in the Supplementary reproduction map.

Fresh C's native wrapper caps are deep 3.0, middle 3.5, and shallow 0.8. The recorded scale request, applied value, and post-scale residual norm are separate fields. The requested C3 profile is $(1.25,2.50,1.25)$; its shallow value is clipped to 0.8 in this run. The Generic Linear and LLH addendum conditions are separately defined and do not equalize realized residual dose. Detailed step- and layer-group records are supplied in the Supplement.

FG-PSNR, FG-LPIPS, FG-SSIM, and Edge-SSIM are computed per object after averaging the six target views. The registered primary comparison is C3--GFL FG-PSNR; secondary and addendum tests follow their locked comparison families as documented in Supplementary Section S2. Fresh C's primary endpoint is foreground-masked PSNR. The historical Table~4 implementation and the Fresh C endpoint are not claimed to be numerically identical; the Fresh C estimate is a new matched held-out result, not a reconstruction of the old $+0.96$~dB contrast. The original 50-object sweep and 24-object probe used separate historical records whose complete current row-level lineage is unavailable, so their implementation details are not inferred from Fresh C.

\paragraph{Statistical reporting}
Fresh C reports the object-level paired mean difference and percentile-bootstrap 95\% interval from 10{,}000 object resamples. The registered primary C3--GFL FG-PSNR contrast is unadjusted as the single primary endpoint. Secondary and addendum families and their multiplicity adjustments are detailed in Supplementary Section S2. CIs outside that prespecified family are pointwise and descriptive. The earlier three-condition 3AFC study is reported separately below as a subjective preference result; it is not a Fresh C endpoint or evidence of texture fidelity.

\subsection{Historical Fixed-Scale Image-Space Diagnostics}
\label{subsec:tradeoff}

Table~\ref{tab:sweep} reproduces the 13-scale sweep reported in the submitted 01549 manuscript. We retain it to preserve the original development question, but the object-level outputs, full run manifest, and metric implementation mapping are not available in the current authority. Its values and Figures~\ref{fig:curve}--\ref{fig:heatmap} are therefore historical diagnostics, not a new reanalysis or confirmatory evidence.

\begin{table}[t]
\centering
\caption{Historical 01549 report of a 13-scale sweep on the stated 50-object set. Values are reproduced from the submission and have not been independently recomputed from row-level outputs; they are retained as development diagnostics.}
\label{tab:sweep}
\small
\begin{tabular}{lcccc}
\toprule
Scale & $\Delta$FG-SSIM$\uparrow$ & POS & $\Delta$Crop-SSIM$\uparrow$ & $\Delta$PSNR$\uparrow$ \\
\midrule
0.75 & $-$0.013 & 20/50 & +0.113 & 7.53 \\
0.90 & +0.024 & 30/50 & +0.137 & 8.32 \\
1.00 & +0.047 & 34/50 & +0.147 & 8.79 \\
1.10 & +0.070 & 36/50 & +0.159 & 9.08 \\
1.15 & +0.071 & 36/50 & +0.157 & 9.25 \\
1.25 & +0.088 & 37/50 & +0.161 & 9.37 \\
1.35 & +0.100 & 39/50 & +0.163 & 9.63 \\
1.50 & +0.113 & 40/50 & +0.163 & 9.76 \\
1.60 & +0.120 & 42/50 & +0.166 & 9.83 \\
1.75 & +0.127 & 43/50 & +0.173 & 9.97 \\
2.00 & +0.134 & 44/50 & +0.169 & 9.96 \\
2.25 & +0.142 & 43/50 & +0.175 & 9.91 \\
2.50 & +0.145 & 44/50 & +0.165 & 9.91 \\
\bottomrule
\end{tabular}
\end{table}

The submitted table reports a rising $\Delta$FG-SSIM over most of the listed scale range and a near-plateau in the crop metric at higher settings. We retain this reported pattern only as historical motivation: absent row-level outputs and a complete run manifest, we do not treat it as an independently confirmed scale-response law. It motivates asking whether a single aggregate structural score fully describes the effect of adapter scaling.

\begin{figure}[t]
\centering
\includegraphics[width=\columnwidth]{fig2.pdf}
\caption{Historical fixed-scale diagnostic from the submitted 01549 paper. The plotted source reports FG-SSIM and Crop-SSIM over adapter scales; the underlying object-level run records are not available for independent reanalysis. It motivates measuring more than one endpoint, but is not a current confirmatory result.}
\label{fig:curve}
\end{figure}

The submitted heatmap depicts object-level variation across the reported scales, but its underlying rows cannot be reconstructed from the current evidence archive. We retain the visual diagnostic as part of the 01549 research history and do not infer a population-level response distribution from it. The updated Fresh C analysis below uses its own frozen cohort and matched inputs.

\begin{figure}[t]
\centering
\includegraphics[width=\columnwidth]{fig3.pdf}
\caption{Historical per-object $\Delta$FG-SSIM heatmap reproduced from the 01549 submission (50 objects, 13 scales). The plot illustrates the originally reported response variation; row-level provenance is incomplete, so it is not used for a current population inference.}
\label{fig:heatmap}
\end{figure}

\subsection{C3 Development Record and Current Evaluation}
\label{subsec:c3}

The submitted paper records 17 uniform, layer-wise, and timestep-conditioned variants on a 24-object development probe and identifies C3 as its chosen profile. This study question remains relevant: does a stage-dependent request offer a useful control setting beyond a single global scale? Because the complete historical row-level data and selection provenance are not available in the current authority, the original Table~\ref{tab:probe} and Figure~\ref{fig:pareto} are retained as historical development records, not as a verified ranking or proof of optimality. Fresh C evaluates the recorded C3 profile without another schedule search.

\begin{table}[t]
\centering
\caption{Representative values reported for the historical 24-object, 17-variant development probe. The full candidate set is shown in Figure~\ref{fig:pareto}. The underlying row-level outputs and selection-time provenance are not available for independent reanalysis; the table records the submitted development result and does not establish an optimum.}
\label{tab:probe}
\footnotesize\setlength{\tabcolsep}{3pt}
\begin{tabular}{llcc}
\toprule
Variant & Reported category & $\Delta$FG-SSIM & $\Delta$Lap Var \\
\midrule
C4 & Late-stage high & 0.105 & $-$0.0120 \\
C2 & Timestep schedule & 0.104 & $-$0.0110 \\
B2 & Layer-wise schedule & 0.090 & $-$0.0135 \\
A\_s2.00 & Uniform & 0.087 & $-$0.0062 \\
C3 & Mid-stage high & 0.084 & $-$0.0006 \\
A\_s2.25 & Uniform & 0.079 & $-$0.0096 \\
A\_s2.50 & Uniform & 0.078 & $-$0.0108 \\
B6 & Layer-wise schedule & 0.000 & $-$0.0027 \\
\bottomrule
\end{tabular}
\end{table}

The original table and plot use shape and texture-variation proxies to describe the development choice. They do not establish that texture was faithful to the reference, nor that C3 was the unique best candidate. In this revision, C3 is therefore described as the historical profile carried forward for matched evaluation, not as a theorem-derived or globally optimal schedule.

The recorded profile uses Equation~\ref{eq:tcas} with $(s_e,s_m,s_l)=(1.25,2.50,1.25)$. This is an inference-time instruction on the existing adapter path; no weights are trained. Its tested value is task- and implementation-specific.

\begin{figure*}[!t]
\centering
\includegraphics[width=0.78\textwidth]{figures/fig7_c3_development.pdf}
\caption{Historical 17-variant development probe from the submitted 01549 paper (24 objects). The axes show changes in FG-SSIM and Laplacian variance; variance is a texture-variation diagnostic, not a fidelity measure. “C3” identifies the selected historical profile, not an optimum or general schedule law. The green region records the historical selection region, not a validated quality threshold. The candidate-level outputs are unavailable for independent reanalysis.}
\label{fig:pareto}
\end{figure*}

The source Figure~5 (26-object texture ratios) is not included as an active result because its object IDs, selection trace, and metric lineage are incomplete. The original 24-object Table~3 is moved to Supplementary Section S10 as a verbatim historical diagnostic, with no texture-fidelity or population inference. Both source assets remain preserved in the authenticated submission and forensic archive; neither decision removes the fixed-scale research question or the development record of C3.

\subsection{Large-Scale Matched Image-Space Evaluation}
\label{subsec:largescale}

\paragraph{Reported main-adapter result in 01549}
The authenticated 01549 submission reports pooled PSNR 18.86 dB for C3 and 17.90 dB for fixed high, a 0.96 dB difference on its original 300-object pool, which includes the 24 development objects. This remains a historical reported result, not a probe-disjoint estimate or a numerical replication by Fresh C. The original blinded preference evidence is retained in Section~\ref{subsec:preference}. The following evaluation uses a different frozen checkpoint and its own paired foreground endpoints.

Fresh C is a revision-era matched image-space study on $N=300$ objects. Its primary run contains No Adapter, global fixed-low (GFL), global fixed-high (GFH), and the historical C3 profile. A separate addendum evaluates Generic Linear and a layer-conditioned low--low--high temporal profile (LLH) on the same 300 object identifiers; it is not an independent cohort. All conditions within a run share the object, reference/target inputs, six target views, initial latents, and seed. Object is the inferential unit, with views averaged within object.

The registered primary contrast is C3 minus GFL on foreground PSNR (FG-PSNR). The point estimate is $+0.52$~dB (95\% percentile-bootstrap CI $[+0.40,+0.64]$); C3 is higher on 206/300 objects (68.7\%), with median paired difference $+0.554514$~dB. Registered secondary endpoints are reported with the multiplicity treatment and full paired statistics in Supplementary Section S2. PSNR, LPIPS, and SSIM quantify image-space agreement on their respective endpoints; they are not direct measures of texture fidelity or 3D material quality.

This primary contrast changes more than temporal placement. Across the 50 steps, the applied deep and middle scale sums are 82.5 each for C3 and 62.5 each for GFL; the shallow sum is 40 for both after the native 0.8 cap. The Fresh C estimate therefore supports the recorded C3 intervention relative to GFL under this wrapper, but cannot be attributed to timing alone or interpreted as a pure schedule-shape effect.

\begin{table*}[t]
\centering
\caption{Fresh C image-space condition means. The primary run contains No Adapter, GFL, GFH, and C3. Generic Linear and LLH are revision-era addendum conditions evaluated on the same object identifiers, not a second independent cohort. Values are means of object-level metrics over six target views.}
\label{tab:freshc-absolute}
\scriptsize\setlength{\tabcolsep}{4pt}
\begin{tabular}{lccccc}
\toprule
Method & $N$ & FG-PSNR$\uparrow$ & FG-LPIPS$\downarrow$ & FG-SSIM$\uparrow$ & Edge-SSIM$\uparrow$ \\
\midrule
\multicolumn{6}{l}{Primary run} \\
No Adapter & 300 & 6.9779 & 0.2288 & 0.4086 & 0.4314 \\
GFL & 300 & 11.3077 & 0.2051 & 0.3500 & 0.4728 \\
GFH & 300 & 12.5777 & 0.1891 & 0.4201 & 0.4961 \\
C3 (TCAS) & 300 & 11.8301 & 0.1998 & 0.3746 & 0.4773 \\
\midrule
\multicolumn{6}{l}{Same-cohort revision-era addendum} \\
Generic Linear & 300 & 12.4230 & 0.1862 & 0.4026 & 0.4880 \\
LLH & 300 & 12.4506 & 0.1876 & 0.3911 & 0.4863 \\
\bottomrule
\end{tabular}
\end{table*}

\begin{table*}[t]
\centering
\caption{Fresh C paired mean differences (A minus B) and pointwise 95\% object-bootstrap CIs. Positive PSNR/SSIM and negative LPIPS favour A. C3--No Adapter/GFH are post hoc; the addendum shares object IDs. Full adjusted families are in Supplement S2.}
\label{tab:freshc-contrasts}
\scriptsize\setlength{\tabcolsep}{3pt}
\begin{tabular}{lcccc}
\toprule
Contrast & FG-PSNR (dB) & FG-LPIPS & FG-SSIM & Edge-SSIM \\
\midrule
C3 $-$ GFL & $+0.522\ [0.404,0.641]$ & $-0.005345\ [-0.006326,-0.004403]$ & $+0.02453\ [0.02120,0.02805]$ & $+0.004494\ [0.003411,0.005586]$ \\
C3 $-$ No Adapter & $+4.852\ [4.553,5.142]$ & $-0.02905\ [-0.03352,-0.02464]$ & $-0.03401\ [-0.04214,-0.02591]$ & $+0.04586\ [0.04114,0.05054]$ \\
C3 $-$ GFH & $-0.748\ [-0.885,-0.609]$ & $+0.01070\ [0.008958,0.012513]$ & $-0.04553\ [-0.04937,-0.04167]$ & $-0.01881\ [-0.02084,-0.01683]$ \\
C3 $-$ Generic Linear & $-0.593\ [-0.837,-0.351]$ & $+0.01351\ [0.01160,0.01549]$ & $-0.02810\ [-0.03197,-0.02439]$ & $-0.01069\ [-0.01300,-0.00850]$ \\
LLH $-$ C3 & $+0.621\ [0.391,0.855]$ & $-0.01214\ [-0.01386,-0.01048]$ & $+0.01658\ [0.01311,0.02003]$ & $+0.009033\ [0.006867,0.011296]$ \\
LLH $-$ Generic Linear & $+0.028\ [-0.015,0.070]$ & $+0.001372\ [0.000762,0.002081]$ & $-0.01152\ [-0.01337,-0.00996]$ & $-0.001660\ [-0.002096,-0.001256]$ \\
\bottomrule
\end{tabular}
\end{table*}

The C3--GFL comparison supports a bounded statement: on this frozen, UID-disjoint cohort and image-space endpoint, the recorded C3 intervention improves mean FG-PSNR relative to GFL. The paired conditions use one frozen checkpoint, common runner, seed policy, inputs, target views, metrics, and recorded native-cap semantics. The historical 01549 v1/v2 checkpoint distinction is preserved, and Fresh C is not presented as a numerical reproduction of the old Table~4 result. This historical difference does not alter the internal pairing or held-out status of Fresh C.

The counter-results delimit that finding. GFH has more favorable cohort means than C3 on all four displayed endpoints. Generic Linear also has more favorable means than C3 on all four; for C3 minus Generic Linear, the mean FG-PSNR difference is $-0.593$~dB (95\% CI $[-0.837,-0.351]$), although the median is $+0.073956$~dB and C3 is higher on 155/300 objects (51.7\%). LLH's FG-PSNR mean exceeds C3 by $+0.621$~dB, while its median difference is $-0.153$~dB and only 140/300 objects favour LLH; its lower mean FG-LPIPS favours LLH. For LLH versus Generic Linear, the FG-PSNR CI includes zero, which is not evidence of equivalence; the FG-LPIPS difference favours Generic Linear. The endpoint means and object-level distributions do not identify a universal winner.

Fresh C has no validated texture-fidelity endpoint. RGB, Laplacian, and gradient responses describe variation: a larger value can represent detail or an artifact. Figure~\ref{fig:c3compare} shows selected full-object/same-ROI fixed-scale and Generic Linear comparisons as qualitative context; they do not estimate prevalence, colorimetric agreement, or texture fidelity.

\begin{figure*}[!t]
\centering
\includegraphics[width=\textwidth]{figures/fig5_freshc_context_roi.pdf}
\caption{Fresh C view 0: (a) Octa\`edre, fixed-scale comparison; (b) Cube 01, Linear counterexample. Columns share the same camera, object, seed and display frame. Red boxes mark the identical pixel crop enlarged below. C3--Linear object-level differences in (b) are $-3.491$ dB FG-PSNR and $+0.08172$ FG-LPIPS. Added patterns differ from the reference despite changes in texture variation. These outcome-selected examples do not estimate prevalence; credits are in Supplementary Table S1.}
\label{fig:c3compare}
\end{figure*}


\subsection{Applied-Scale Budget and Temporal Placement}
\label{subsec:temporal-placement}

Fresh C compares the recorded C3 profile with GFL under the same frozen checkpoint, runner, seed policy, inputs, and target views. Its applied deep/middle scale sums are 82.5 for C3 and 62.5 for GFL, while the shallow sum is 40 for both after the 0.8 cap. The registered C3--GFL improvement therefore combines a change in middle-third strength with a larger per-group nominal budget; it does not isolate temporal placement.

A separate post-inspection placement analysis used the clean-v2 strict-276 cohort and a common runner within that run. It is not the Fresh C cohort, and it is not an independent replication of the historical 24-object selection probe. Temporal L--L--H means low requests during early and middle steps and high requests during late steps; it uses a different shallow-group control from Fresh C's layer-conditioned LLH addendum. The latter requests (1.25,1.25,2.50) in deep/middle and (0.50,0.50,0.75) in shallow; the older global placement comparison keeps shallow capped at 0.8. Neither is identical to the new 16-step Late-high condition. With 50 steps, the strict-276 run assigns 17/16/17 steps to the three windows. C3 and temporal L--L--H have near-matched applied-scale sums (82.50 versus 83.75 in deep and middle; 40 in shallow), but their post-scale residual magnitudes are not matched.

\begin{table*}[t]
\centering
\scriptsize
\setlength{\tabcolsep}{4pt}
\caption{Post-inspection strict-276 placement analysis. Entries are C3 minus temporal L--L--H with pointwise 95\% object-bootstrap intervals; positive PSNR/SSIM and negative LPIPS differences favor C3. Temporal L--L--H is favored on all six formal endpoints. The analysis is separate from Fresh C, and its scalar sums do not match realized residual dose.}
\label{tab:temporal-placement}
\begin{tabular}{lcc}
\toprule
Endpoint & Mean difference [95\% CI] & Direction \\
\midrule
Full-view PSNR (dB) & $-0.268100$ [$-0.311101,-0.227797$] & L--L--H \\
FG-PSNR (dB) & $-0.177084$ [$-0.230539,-0.130784$] & L--L--H \\
FG-SSIM & $-0.008851$ [$-0.011455,-0.006291$] & L--L--H \\
Edge-SSIM & $-0.035301$ [$-0.036581,-0.034031$] & L--L--H \\
FG-LPIPS & $+0.001896$ [$+0.001205,+0.002578$] & L--L--H \\
Full-view LPIPS & $+0.008658$ [$+0.007465,+0.009831$] & L--L--H \\
\bottomrule
\end{tabular}
\end{table*}

Against the fixed-mean profile, C3's strict-276 FG-PSNR difference is $+0.502875$~dB (95\% CI $[+0.462018,+0.542244]$). These results show sensitivity to the tested stage allocation, not a universally preferred placement or an equal-realized-dose effect.

Fresh B and Core-7 remain separate retrospective sensitivity analyses in Supplementary Sections S3, S5 and S21; they are not pooled with Fresh C or used to replicate the historical headline.

\subsection{Matched Stage Allocation at Equal Applied Budget}
\label{subsec:matched-budget300}
The six-condition matrix contains 1,800 object--condition rows, with complete input-hash and common-seed checks. It is a post-inspection addendum on Fresh C's 300 objects, not an independent holdout. Early-, middle-, and late-high allocate the same 16 high steps to updates 0--15, 17--32, and 34--49. Deep and middle use low 1.25 and high 2.50; shallow remains capped at 0.8. Fixed mean uses 1.65 before the same cap. Thus all four allocation conditions have per-group applied sums (82.5, 82.5, 40), or 205 in total. Fixed Low and Fixed High are separate strength controls and do not share that budget. Generic Linear remains a strong practical baseline in Table~\ref{tab:freshc-absolute}, rather than being silently treated as equal-budget.

\begin{table*}[t]
\centering\scriptsize\setlength{\tabcolsep}{4pt}
\caption{Fully matched Fresh C control means ($N=300$ per condition). Timing and fixed-mean rows share exactly the same per-layer applied-scale sums. This same-cohort addendum is post-inspection.}
\label{tab:matched-budget300}
\begin{tabular}{lrrrrr}
\toprule
Control & FG-PSNR$\uparrow$ & FG-LPIPS$\downarrow$ & FG-SSIM$\uparrow$ & Edge-SSIM$\uparrow$ & PSNR winners \\
\midrule
GFH / Fixed High & 12.58 & 0.1891 & 0.4201 & 0.4961 & 172 \\
GFL / Fixed Low & 11.31 & 0.2051 & 0.3500 & 0.4728 & 37 \\
Fixed Mean & 11.96 & 0.1978 & 0.3864 & 0.4843 & 2 \\
Early-high & 11.33 & 0.2045 & 0.3696 & 0.4765 & 25 \\
C3 / Middle-high & 11.83 & 0.1998 & 0.3746 & 0.4773 & 0 \\
Late-high & 12.39 & 0.1951 & 0.3788 & 0.4914 & 64 \\
\bottomrule
\end{tabular}
\end{table*}

C3 minus Fixed Mean is $-0.13$ dB (pointwise 95\% object-bootstrap CI $[-0.17,-0.10]$). C3 minus Late-high is $-0.56$ dB (pointwise 95\% object-bootstrap CI $[-0.63,-0.49]$). C3 minus Early-high is $+0.50$ dB (pointwise 95\% object-bootstrap CI $[+0.37,+0.62]$). C3 also trails Fixed Mean and Late-high on FG-LPIPS, FG-SSIM and Edge-SSIM. Complete paired endpoint distributions, 50,000 object-bootstrap resamples and post hoc Holm-adjusted contrasts are supplied in Supplementary Section S16. The original stage-centered contrast, C3 minus the mean of Early-high and Late-high FG-PSNR, is $-0.034$ dB (95\% CI $[-0.117,+0.051]$; paired sign-flip $p=0.4324$). It likewise does not support a middle-window advantage; the later pairwise analysis is explicitly post-inspection. The rank count is specific to FG-PSNR and these six conditions; it does not establish a universal quality winner. We retain C3's development identity but reject a generally optimal middle window. Residual amplitudes differ along the resulting trajectories: this restricts a further mechanism claim, while the equal-budget allocation comparison still measures the total effect of stage allocation.

\subsection{Layer-by-Window Sensitivity}
\label{subsec:layerwindow}

The A2 and A3b maps test response heterogeneity over three adapter groups and five denoising windows. Their object-cluster statistics describe variation within the tested interventions; they are not effect estimates for a dose-independent schedule mechanism.

\begin{table}[t]
\centering
\scriptsize
\setlength{\tabcolsep}{3pt}
\caption{Object-cluster tests of the observed layer-by-window response maps. A2 is exploratory; A3b is retrospective Fresh B sensitivity because the preserved history cannot authenticate pre-run availability of its lock. Interaction share and cell RMSE summarize the observed map, not a causal effect.}
\label{tab:layerwindow-main}
\begin{tabular}{llrrr}
\toprule
Map, cohort & Endpoint & Wald $W(8)$ & Reported $p$ & Share; cell RMSE \\
\midrule
A2, $N=300$ & FG-PSNR & 1879.963 & $<10^{-300}$ & 56.8\%; 0.1222~dB \\
A2, $N=300$ & FG-LPIPS & 481.081 & $8.04\times10^{-99}$ & 43.5\%; 0.001695 \\
A3b, $N=150$ & FG-PSNR & 525.700 & $3.43\times10^{-107}$ & 36.0\%; 0.0205~dB \\
A3b, $N=150$ & FG-LPIPS & 225.015 & $5.37\times10^{-43}$ & 14.0\%; 0.000342 \\
\bottomrule
\end{tabular}
\end{table}

The response maps describe heterogeneity under the tested controls. Cross-layer raw norms and retrospective complexity associations do not establish an optimal schedule or a causal mechanism; full diagnostics are in Supplementary Sections S3--S4, S12 and S21.

\subsection{Appearance and Color Failure Analysis}
\label{subsec:appearance-color}

The selected color diagnosis contains 26 failure-enriched objects: six earlier figure cases and 20 Fresh C panels. It is not a random sample, a new holdout, or a population estimate. No Adapter, GFL, and LLH were compared on the same inputs, checkpoint, and runner with six target views per object (78 object-condition records). The following paired object-bootstrap intervals are descriptive (20,000 resamples; seed 61008), and no confirmatory multiplicity family was defined.

\begin{figure*}[t]
\centering
\includegraphics[width=\textwidth]{figures/fig6_color_diagnostic_freshc_context_roi.pdf}
\caption{Fresh C view 0, E5-5-2. The same frame shows reference, disabled adapter, Fixed Low and layer LLH. Geometry conditioning changes alignment and patterned appearance, but the brown texture does not reproduce the gray reference. This previously viewed example was selected by an existing LPIPS stratum, not for color. It is not a reconstruction of the legacy purple-shift panels; attribution is in Supplementary Table S1.}
\label{fig:color-failure}
\end{figure*}

\begin{table}[t]
\centering
\scriptsize
\setlength{\tabcolsep}{3pt}
\caption{Selected paired results from the failure-enriched 26-object color diagnosis. Negative CIEDE2000 differences indicate lower reference-color error for the first condition. The set is descriptive and does not estimate population prevalence.}
\label{tab:color-summary}
\begin{tabular}{p{2.3cm}p{3.3cm}p{1.3cm}}
\toprule
Contrast & Mean CIEDE2000 difference [95\% CI] & Objects \\
\midrule
GFL -- No Adapter & $-19.654$ [$-26.370,-13.248$] & 24/26 lower \\
LLH -- GFL & $+3.583$ [$+0.462,+6.712$] & -- \\
\bottomrule
\end{tabular}
\end{table}

In the selected 26-object diagnosis, GFL reduces mean reference-color error relative to No Adapter by 19.65 CIEDE2000 units (95\% paired interval $[-26.37,-13.25]$; 24/26 favorable). LLH increases color error relative to GFL by 3.58 while reducing local variation. These paired results distinguish conditioning response from reference fidelity; the set is failure-enriched and does not estimate prevalence. Figure~\ref{fig:color-failure} shows a traced Fresh C example, not a reconstruction of the two unmapped legacy purple-shift panels. No repair is claimed.

Historical bake cohorts, execution-stack checks and unsuccessful repair candidates are documented in Supplementary Sections S13, S18--S19 and S21. They are supporting diagnostics rather than additional evidence for C3 quality. The new calibrated bake pilot is evaluated separately below.

\paragraph{Fresh C full-object and same-ROI examples.}
\label{subsec:visual}

The submitted 26-object texture audit and its quantitative Figure~5 cannot be re-traced to object IDs, selection criteria, raw values, and an executable metric definition in the current evidence archive. We therefore remove its 92\% prevalence statement, its “true shape gain” classification, and its RGB-standard-deviation/gradient ratios from the scientific claims. Those statistics would describe variation, not texture fidelity, even if their source records were complete.

Figure~\ref{fig:qualitative} presents two Fresh C examples selected where C3 improves over Fixed Low on FG-PSNR, FG-LPIPS and FG-SSIM. The caption reports their object-mean values across six views; the 300-object comparison in Table~\ref{tab:freshc-absolute} gives cohort-level results. These metric-selected examples illustrate the measured direction and do not estimate its frequency. Figure~\ref{fig:c3compare} shows a fixed-scale comparison and a Generic Linear example with the same-pixel ROI enlarged below; Supplementary Figure~S1 retains the six-view Boat comparison.

\begin{figure*}[!t]
\centering
\includegraphics[width=\textwidth]{figures/fig4_metric_favorable.pdf}
\caption{Metric-selected Fresh C examples comparing C3 with Fixed Low: (a) Building 03 and (b) Cube 01. Across six generated views, C3--Fixed Low object-mean differences are FG-PSNR $+2.64$ and $+1.79$ dB, FG-LPIPS $-0.002$ and $-0.024$, FG-SSIM $+0.034$ and $+0.185$, and Edge-SSIM $+0.013$ and $-0.003$, respectively. Each panel shows the corresponding view with a shared crop across methods. Higher FG-PSNR/SSIM and lower FG-LPIPS indicate favorable values. Both cases are drawn from the 178/300 objects favorable on all three metrics; the full cohort results are in Table~\ref{tab:freshc-absolute}. Credits are in Supplementary Table S1.}
\label{fig:qualitative}
\end{figure*}




\subsection{Calibrated Baked Appearance on New Objects}
\label{subsec:new-bake}
We froze ten previously unbaked Fresh C objects in UID order using input availability, native UV/material compatibility, triangle-count and camera-mask alignment criteria. Of the 300 screened inputs, only ten passed all criteria before the cutoff; the planned maximum was twenty. The ten are not new to the GeoTex image domain, but their baked outcomes had not been inspected during selection. This is a compatibility-selected pilot, not a representative twenty-object confirmation. Five earlier bake objects are reported separately.

Seven conditions share each object's source mesh, reference, seed, six projection cameras, masks, 256-pixel source resolution, atlas resolution, hole filling and eleven unseen evaluation cameras. Original geometry, UVs and legal base-color material semantics are retained and checked. Unseen cameras do not enter the bake. GT-256 uses the same source resolution, masks and bake parameters and has exactly the same raw atlas coverage as every generated condition. Native REPEAT sampling passes boundary and integer-period tests; pole-view RGB orientation is corrected without rotating raw-camera masks. These calibrations control projection artifacts; evaluation is unlit base color, not PBR fidelity. Supplementary Section~S18 records the checks and object-level data. The Constant 1.0 control is capped at the shallow group and is not the unset-scale GeoTex default. No generated output for that unmodified path is available in this comparison, so the corresponding reviewer request remains partially unresolved.

\begin{table}[ht]
\centering\scriptsize
\caption{New ten-object unlit unseen-view bake pilot. Eleven views are averaged per object. Higher PSNR/SSIM and lower LPIPS/CIEDE2000 are favorable. Constant 1.0 is an explicitly capped request, not the unset-scale pipeline. The GT row is a projection calibration control; all seven generated conditions and all ten eligible objects are retained.}
\label{tab:new-bake}
\begin{tabular}{lrrrr}
\toprule Condition & FG-PSNR & FG-LPIPS & FG-SSIM & CIEDE2000 \\
\midrule
GT-256 bake (calibration) & 22.38 & 0.0590 & 0.7683 & 4.44 \\
Constant 1.0 (capped) & 11.93 & 0.1354 & 0.3504 & 24.18 \\
Fixed Low & 12.48 & 0.1317 & 0.3835 & 22.71 \\
Fixed High & 12.29 & 0.1305 & 0.4587 & 23.27 \\
Fixed Mean & 12.62 & 0.1303 & 0.4193 & 22.32 \\
C3 / Middle-high & 12.48 & 0.1302 & 0.4059 & 22.16 \\
Late-high & 12.92 & 0.1277 & 0.4259 & 21.65 \\
Generic Linear & 11.60 & 0.1276 & 0.4458 & 24.71 \\
\bottomrule\end{tabular}
\end{table}

The internally predeclared primary contrast was Late-high minus Fixed Mean on unseen-view FG-LPIPS, proposed from the earlier pilot and tested on the ten newly selected bake objects. Its mean is $-0.00265$ (95\% object-bootstrap CI $[-0.00414,-0.00137]$; exact paired sign-flip $p=0.00195$), with all ten objects favorable. This is approximately a 2.0\% reduction relative to Fixed Mean's mean LPIPS, and supports a small relative metric improvement on this interface, without establishing a perceptible or usable texture improvement. Generic Linear has a slightly lower mean LPIPS but lower PSNR; Late-high does not dominate every endpoint or control.

C3--Fixed High has mean PSNR $+0.19$~dB (95\% CI $[-0.65,+0.87]$) and LPIPS $-0.00036$ ($[-0.00396,+0.00336]$); neither separates under the four-comparison secondary Holm family (both adjusted $p=1$). Its SSIM is lower by 0.05274 ($[-0.07459,-0.03170]$; Holm $p=0.00781$). Thus the old five-object C3 PSNR advantage does not repeat stably on this new pilot. We retain C3 as the historical candidate and do not rename Late-high as C3.

Across the seven conditions, mean within-object Spearman correlations between generated-view and baked unseen-view rankings are 0.846 for PSNR, 0.804 for LPIPS, 0.746 for SSIM and 0.907 for Edge-SSIM. Broad rankings are positively associated, while the winning condition changes on 7/10 objects for LPIPS (1/10 for PSNR). These are different statements: we do not claim a general reversal of rankings. Full-object Figure~\ref{fig:baked-context} also shows persistent appearance errors, including the Tire's brown output against a gray reference. A lower LPIPS does not establish faithful material reproduction. Across the displayed objects, similar brown patterns replace distinct source colors, facade structure, wheel detail and printed markings. Source and GT-256 bakes retain much of that appearance, so projection alone does not explain these failures. Their common cause is unresolved; the evidence establishes relative responses of the frozen deployment rather than usable texture recovery.

\begin{figure*}[t]
\centering
\includegraphics[width=\textwidth]{figures/fig8_baked_unseen_context.pdf}
\caption{Unseen camera 004: (a) LG Remote, (b) Tire. Source and GT-256 bake calibrate projection; Fixed Mean and Late-high isolate the principal comparison. All columns share one display frame, without color/detail editing. GT retains much of the source appearance; generated results retain brown-pattern errors despite a small cohort-level LPIPS improvement. Complete seven-condition comparisons are in Supplementary Figures S2--S8. Credits: mohd.akmal.mahyudin and Koto3D Stephane Vogley, CC Attribution (Supplement S18).}
\label{fig:baked-context}
\end{figure*}

\subsection{Cross-Interface Evaluation}
\label{subsec:cross-interface}

Table~\ref{tab:cross-interface} places the primary GeoTex path beside the auxiliary MV-Adapter and MVDiffusion interventions. The methods use different conditioning interfaces and scale semantics; object-level results are not pooled across systems.

\begin{table*}[t]
\centering
\scriptsize
\setlength{\tabcolsep}{3pt}
\caption{Interface-specific intervention evidence. Nominal budgets are described within each implementation; no cross-interface realized-dose match is available. These rows are not architecture-matched replications. The MVDiffusion contrast is reported as raw condition A minus condition B; lower FG-LPIPS, CIEDE2000, and GT-relative texture error are favorable, whereas higher PSNR, FG-SSIM, and Edge-SSIM are favorable.}
\label{tab:cross-interface}
\begin{tabular}{p{2.4cm}p{4.0cm}p{4.2cm}p{4.5cm}}
\toprule
System and injected signal & Control axis and budget & Cohort and principal result & Evidence grade and limitation \\
\midrule
GeoTex-Adapter, Fresh C &
Five-channel normal/depth/foreground adapter residuals added to UNet ResNet features; temporal requests with native layer caps. C3--GFL applied deep/middle sums differ (82.5 vs.\ 62.5); shallow sums are both 40. &
$N=300$. C3--GFL FG-PSNR is $+0.522$~dB [0.404, 0.641], but GFH and Generic Linear have more favorable means than C3 on all four displayed endpoints. &
Matched, same-runner Fresh C image-space comparison; identifies the named paired contrast, not timing alone, a schedule optimum, or baked material quality. \\
MV-Adapter, layer/window probe &
Condition guider plus decoupled multi-view attention; one of four down-block points is raised from 0.75 to 1.00 in one of five 10-step windows. Group residual-dose ratio is 1:0.57:1.46; no equal-realized-dose map was run. &
98/99 planned objects. The exploratory four-metric Holm-corrected cluster-Wald family found no interaction detected after correction (primary FG-LPIPS Holm $p=0.303$). &
Interface-specific layer/window sensitivity; not C3 temporal transfer. A corrected non-detection does not show no interaction or equivalence. \\
MV-Adapter, layer/time profile panel &
Four condition-guider features were scaled per point; L-FIX, L-LHL, and L-LLH add 228 outputs to the existing global schedules on the same 76-object Exact holdout. &
Layer-wise/global contrasts are metric-dependent. The two within-profile temporal contrasts show no consistent pattern; their 30 pointwise intervals are unadjusted. &
Retrospective: protocol and outputs first appear in one preserved commit. This feature-scaling interface differs from GeoTex residual correction and does not establish TCAS transfer. \\
MVDiffusion, compatibility interface &
Correspondence-aware decoder-side CPBlocks; the tested multiplier is a product of a layer factor and temporal scale over five blocks. The 75-object panel used the compatibility base; a later check found 1,552/1,552 loaded tensors and 12/12 views identical to the native mirror on one object. &
$N=75$, 12-view output grids, six frozen target views scored per object, and 50 steps. For raw L--LLH minus G--FL means: PSNR $-0.4538$, FG-SSIM $-0.01939$, Edge-SSIM $-0.005941$, FG-LPIPS $+0.002745$, CIEDE2000 $+1.3730$, and GT-relative texture error $-0.08229$. Full pointwise 95\% CIs for all endpoints and five contrasts are in Supplementary Table~S12. &
Within-interface, mixed endpoint comparison on a different conditioning path; the 75-object panel itself was not rerun on the mirror. No shared scale semantics, cross-system dose match, or architecture-causal inference. \\
\bottomrule
\end{tabular}
\end{table*}

Under the raw A-minus-B convention, L--LLH has higher CIEDE2000 (worse reference-color agreement) than G--FL, lower PSNR/FG-SSIM/Edge-SSIM, lower GT-relative error against the recorded variation statistics, and no separated FG-LPIPS difference. The lower variation-statistic error is not evidence of texture fidelity. The earlier negative CIEDE2000 value was an improvement-coded effect, not a raw contrast; the direction audit and all five comparisons are reported in Supplementary Table~S12.

For the 98-object MV-Adapter probe, the original additive-residual bootstrap had an invalid zero-power null and was not used for interaction inference. The exploratory object-cluster Wald test is reported with its corrected family; its non-detection does not establish absence of interaction. In the separate 76-object Exact-holdout panel, layer-wise versus global contrasts trade PSNR, color, and ground-truth-relative texture endpoints against FG- and Edge-SSIM. For temporal placement, only one of six pointwise 95\% intervals excludes zero in L-LLH versus L-FIX (direction-adjusted CIEDE2000 $+0.072790$ [$+0.004371,+0.156696$]); only one does so in G-LLH versus G-LHL (Edge-SSIM $-0.000261$ [$-0.000522,-0.000032$]). The 30 intervals are descriptive and unadjusted. The protocol and outputs first appear in the same preserved commit, so this is retrospective sensitivity evidence, not confirmation or an equivalence test. The MVDiffusion panel used a compatibility deployment; a separate check matched its loaded tensors and one object's rendered views to a native mirror, but the full panel was not rerun. These probes remain interface-specific and do not establish architecture-invariant scheduling or transfer of the GeoTex residual intervention.


\paragraph{New native MV-Adapter temporal comparison.}
\label{subsec:new-mva}
We froze fifty Fresh C inputs not previously inspected through MV-Adapter, excluding all recorded earlier native and profile-panel objects. The input-only ordered selection used availability and native material/geometry compatibility. The native mesh normalization, six orthographic cameras and nvdiffrast rasterizer also generate evaluation targets; an independent source-UV sampling check validates this target path. Three textured new objects pass official/deployed and constant-default-path identity checks: 18 RGB images, conditioning features and injected residuals are bitwise identical. The old seven-object quality archive fails camera calibration and is superseded, not pooled.

The seven conditions share object inputs, seed, geometry and evaluation code. With 50 steps, low/high are 0.75/1.0; each stage-high schedule has 16 high steps, Fixed Mean is 0.83, and Generic Linear is 0.75--0.91. These five mean-matched controls each sum to 41.5 per scaled feature, within this interface. The fifty-object matrix is complete (350 rows), including thirteen texture-absent sources with white-albedo targets and a documented loader compatibility path. The other 37 textured objects are a descriptive subset, not a replacement primary cohort. Targets are unlit native-camera albedo, not Fresh C's differently normalized reference renders.

\begin{table}[ht]\centering\scriptsize
\caption{Native-camera MV-Adapter comparison, fifty input-selected objects. Views are averaged within object. Scale values and endpoints are specific to this interface, and are not numerical doses shared with GeoTex.}\label{tab:new-mva}
\begin{tabular}{lrrrr}
\toprule Condition & FG-PSNR & FG-LPIPS & FG-SSIM & Edge-SSIM \\
\midrule
Native / Fixed High & 12.60 & 0.1280 & 0.5041 & 0.3526 \\
Fixed Low & 12.73 & 0.1270 & 0.5066 & 0.3475 \\
Fixed Mean & 12.68 & 0.1279 & 0.5068 & 0.3509 \\
Early-high & 12.60 & 0.1296 & 0.5041 & 0.3503 \\
Middle-high & 12.71 & 0.1262 & 0.5061 & 0.3470 \\
Late-high & 12.74 & 0.1260 & 0.5075 & 0.3514 \\
Generic Linear & 12.72 & 0.1263 & 0.5072 & 0.3501 \\
\bottomrule\end{tabular}\end{table}

The internally predeclared primary, Middle-high minus Fixed Mean FG-PSNR, is $+0.029$~dB (95\% object-bootstrap CI $[-0.019,+0.081]$; paired Monte Carlo sign-flip $p=0.281$; 29/50 favorable). This does not support a stable middle-window gain or establish equivalence. The remaining 104 pair--metric tests use one Holm family. Middle-high has lower Edge-SSIM than Fixed Mean (adjusted $p=0.00104$); some other placement contrasts separate, including Late-high versus Early-high PSNR. Such endpoint-dependent differences do not demonstrate transferable C3 optimality. Complete contrasts, correction rules and the textured subset are in Supplementary Section~S17. This is a native-interface validation of the scheduling question, not architecture-matched replication of GeoTex's capped residual intervention.

\begin{figure*}[t]
\centering
\includegraphics[width=\textwidth]{figures/fig9_mvadapter_native.pdf}
\caption{MV-Adapter view 0: (a) LG Remote, (b) Tire. Target, Fixed Mean, Middle-high and Late-high share the native camera and recorded object seed. Original backgrounds are retained, with the same crop across conditions. These first two input-order, rights-cleared illustrations show appearance differences without establishing a stable middle-window gain. Figure S9 includes all seven controls; attribution is in Supplement S18.}
\label{fig:native-mva-visual}
\end{figure*}

\subsection{First-round cluster-aware preference study}
\label{subsec:preference}

The 01549 first-round revision reported a blinded three-alternative forced-choice study on 30 objects with 24 participants. Participants compared conservative fixed-scale, aggressive fixed-scale, and C3 outputs for texture naturalness, shape consistency, and overall quality; condition order was randomized per trial. Each criterion had 720 first-choice decisions. Table~\ref{tab:preference} reports the aggregate counts; participant-level responses are not redistributed with this revision.

\begin{table}[t]
\centering\scriptsize
\caption{First-round 01549 3AFC results retained as reported. Percentages use 720 decisions per criterion; preference is not a measure of reference texture fidelity.}
\label{tab:preference}
\setlength{\tabcolsep}{3pt}
\begin{tabular}{lccc}
\toprule
Method & Texture naturalness & Shape consistency & Overall quality \\
\midrule
$s=1.25$ & 332/720 (46.1\%) & 96/720 (13.3\%) & 176/720 (24.4\%) \\
$s=2.50$ & 104/720 (14.4\%) & 331/720 (46.0\%) & 126/720 (17.5\%) \\
C3 (TCAS) & 284/720 (39.4\%) & 293/720 (40.7\%) & 418/720 (58.1\%) \\
\bottomrule
\end{tabular}
\end{table}

The reported analysis used a two-level participant/object bootstrap with 10,000 resamples. C3's overall-quality share was 58.1\% (95\% CI [50.1\%, 65.8\%]); its differences from the conservative and aggressive baselines were +33.6 points (CI [19.3, 47.2]) and +40.6 points (CI [29.4, 51.8]). The single-criterion advantages for texture naturalness and shape consistency had intervals crossing zero. This bounded preference result does not validate texture fidelity, unseen-view quality, or the Fresh C cohort.

\subsection{FAC controlled re-examination reported in 01549}
\label{subsec:facreported}

A controlled re-examination of Full Adaptive Correction (FAC) was previously reported, with the base model, VAE, and GeoTex-Adapter frozen and roughly $6\times10^3$ controller parameters trainable. The stated protocol used the enhanced latent denoising objective (foreground weight 1.7, edge weight 1.3, and a latent-SSIM term), AdamW at peak learning rate $10^{-4}$ with 100 warm-up steps, and 2,000 training steps on a training pool reported as 1,706 objects and disjoint from evaluation. It reported a 300-object paired evaluation: every trained configuration reduced FG-PSNR versus TCAS by 1.10--2.94 dB, with at most 48/300 wins; the untrained C3-envelope control on 12 objects had $\Delta$FG-PSNR $=-0.18$ dB (95\% CI [$-0.42,+0.05$]). The seven reported configurations and endpoint values are retained in Supplementary Table~S16.

These FAC outcomes are retained as previously reported results. They are not presented as a new 1008 training run. The current archive audit found that a recovered configuration points to a 1,200-object list rather than the 1,706-object pool stated in 01549, and does not bind the training seed and run outputs in one manifest. We therefore disclose the remaining reproducibility gap and make no code-release claim; TCAS remains the sole training-free method.

\section{Discussion}
TCAS turns an existing geometry-adapter multiplier into an explicit temporal control without additional training. The original 01549 preference study favors C3 overall, and Fresh C supports a paired FG-PSNR gain over Fixed Low. These are useful but distinct findings: the preference concerns the original stimuli, while Fresh C tests a retrained checkpoint and foreground image agreement. Together they motivate evaluating the interaction between strength, stage allocation and the endpoint of interest rather than relying on one structural metric.

The strong controls make that evaluation informative. Fixed High and Generic Linear exceed C3 on the displayed Fresh C means; four matched-budget controls favor Fixed Mean and Late-high over the recorded middle window. Equal applied-scale sums permit comparison of allocation policies even though evolving adapter activations prevent a residual-dose-independent mechanism interpretation. This is an implementation-specific response, not a reason to discard temporal control or infer that C3 is generally optimal.

The downstream pilot supplies a small Late-high--Fixed Mean LPIPS benefit, with positive broad image/bake rank association and a different LPIPS winner on seven of ten objects. GT-256 calibration and Figures~5--8 also show substantial appearance discrepancies, so the relative metric gain does not establish reference-faithful material recovery. Native MV-Adapter testing adds a valid second-interface comparison whose primary interval crosses zero. The combined evidence supports reporting strength controls, paired image outcomes and projected appearance together, with the useful allocation tied to the tested interface.

\section{Limitations}
\label{sec:limitations}
Historical 01549, Fresh C, native MV-Adapter and the bake pilots retain separate checkpoint and cohort identities. Fresh C does not directly replicate the original pooled 0.96 dB estimate; its allocation addendum is post-inspection on the same 300 objects. The original preference study is not a current-checkpoint or baked-texture preference test. Historical records with incomplete row-level provenance remain supporting reports (Supplementary Sections S8--S11 and S19).

The new bake pilot contains ten technically compatible objects and evaluates unlit base color. Actual appearance failures persist; the true unset-scale GeoTex output, validated seams and PBR/relighting are absent. Constant 1.0 (capped) is not that default. MV-Adapter scales a different conditioning signal, so numerical strengths and preferred windows do not transfer automatically. Texture-variation proxies are not fidelity measures, and no equivalence, non-inferiority, general optimality or independently derived CAI claim is made. FAC full replay remains limited by missing training records and weights. Detailed diagnostics and recovery boundaries are retained in Supplementary Sections S13 and S18--S21.

\section{Conclusion}
\label{sec:conclusion}
TCAS provides a training-free stage-dependent control for geometry-adapter strength in multi-view texture generation. The C3 profile and original development and preference studies retain their historical identity. A matched Fresh C evaluation supports C3 over Fixed Low, while strong fixed/linear controls and equal-budget allocation comparisons delimit that benefit. Calibrated unseen-view baking supplies a small named Late-high LPIPS result, and native MV-Adapter testing evaluates the same stage question on a second interface without a separated primary gain. These findings establish a concrete, reproducible control-and-evaluation study for the tested texturing system; they do not imply a universal middle-stage optimum or resolved reference-appearance fidelity.

\section*{Data availability}

Object-level revision metrics, frozen protocols/configuration records, verified inference and analysis scripts, and the LaTeX source required to rebuild the manuscript and Supplement are included in the author-review package. The standalone reproduction archive supports statistic recomputation without model weights; fresh inference additionally requires the recorded weights, runtime and licensed source assets. Raw generated-image archives, trained checkpoints, source 3D assets whose redistribution rights are not established, and participant-level responses are not included. Additional analysis records may be requested from the corresponding author, subject to licensing, privacy, and archive constraints.

\section*{Declaration of generative AI and AI-assisted technologies in manuscript preparation}
OpenAI Codex assisted with revision drafting, implementation and reproducible statistical and figure code. Figure~\ref{fig:overview} is a code-generated vector schematic; the earlier generated conceptual artwork is excluded from this revision. Actual input/output thumbnails and experimental displays come from the recorded research runs and were not retouched with an image-generation tool. Numerical tables are computed from delivered object-level records. Final scientific review, accuracy and responsibility remain with the authors.

\section*{Declaration of competing interests}

The authors declare that they have no known competing financial interests or personal relationships that could have appeared to influence the work reported in this paper.


\begingroup
\linespread{0.94}\selectfont
\begin{thebibliography}{45}
\bibitem{liu2023zero1to3}
Liu R, Wu R, Van Hoorick B, et al. Zero-1-to-3: Zero-shot one image to 3D object. \textit{Proceedings of the IEEE/CVF International Conference on Computer Vision}, 2023.

\bibitem{shi2023zero123pp}
Shi R, Chen H, Zhang Z, et al. Zero123++: A single image to consistent multi-view diffusion base model. arXiv:2310.15110, 2023.

\bibitem{liu2023syncdreamer}
Liu Y, Lin C, Zeng Z, et al. SyncDreamer: Generating multiview-consistent images from a single-view image. \textit{International Conference on Learning Representations}, 2024.

\bibitem{long2023wonder3d}
Long X, Guo Y C, Lin C, et al. Wonder3D: Single image to 3D using cross-domain diffusion. \textit{Proceedings of the IEEE/CVF Conference on Computer Vision and Pattern Recognition}, 2024.

\bibitem{shi2024mvdream}
Shi Y, Wang P, Ye J, et al. MVDream: Multi-view diffusion for 3D generation. \textit{International Conference on Learning Representations}, 2024.

\bibitem{wang2023imagedream}
Wang P, et al. ImageDream: Image-prompt multi-view diffusion for 3D generation. arXiv:2312.02201, 2023.

\bibitem{xu2024instantmesh}
Xu J, et al. InstantMesh: Efficient 3D mesh generation from a single image with sparse-view large reconstruction models. \textit{Advances in Neural Information Processing Systems}, 2024.

\bibitem{ke2024era3d}
Li P, Liu Y, Long X, et al. Era3D: High-resolution multiview diffusion using efficient row-wise attention. \textit{Advances in Neural Information Processing Systems}, 2024.

\bibitem{chen2023text2tex}
Chen D Z, et al. Text2Tex: Text-driven texture synthesis via diffusion models. \textit{Proceedings of the IEEE/CVF International Conference on Computer Vision}, 2023.

\bibitem{richardson2023texture}
Richardson E, Metzer G, Alaluf Y, et al. TEXTure: Text-guided texturing of 3D shapes. \textit{ACM SIGGRAPH 2023 Conference Proceedings}, 2023.

\bibitem{poole2023dreamfusion}
Poole B, Jain A, Barron J T, Mildenhall B. DreamFusion: Text-to-3D using 2D diffusion. \textit{International Conference on Learning Representations}, 2023.

\bibitem{chen2023fantasia3d}
Chen R, Chen Y, Jiao N, Jia K. Fantasia3D: Disentangling geometry and appearance for high-quality text-to-3D content creation. \textit{Proceedings of the IEEE/CVF International Conference on Computer Vision}, 2023.

\bibitem{shao2025mvpainter}
Shao M, Xiong F, Sun Z, Xu M. MVPainter: Accurate and detailed 3D texture generation via multi-view diffusion with geometric control. arXiv:2505.12635, 2025.

\bibitem{kynkanniemi2024limited}
Kynk\"a\"anniemi T, Aittala M, Karras T, et al. Applying guidance in a limited interval improves sample and distribution quality in diffusion models. \textit{Advances in Neural Information Processing Systems}, 2024.

\bibitem{perla2026survey}
Perla S R K, Zhang H, Mahdavi-Amiri A. Advances in neural 3D mesh texturing: A survey. \textit{Computer Graphics Forum}, 2026, e70392.

\bibitem{guzman2026adaptive}
Guzm\'an J E, Mould D, Paquette E. Adaptive multiresolution exemplar-based texture synthesis on animated fluids. \textit{Computers \& Graphics}, 2026, 134: 104490.

\bibitem{lv2025interpretable}
Lv X, Wu Z, Xu J, et al. Interpretable procedural material graph generation via diffusion models from reference images. \textit{The Visual Computer}, 2025, 41(13): 11195--11205.

\bibitem{zeng2024paint3d}
Zeng X, Chen X, Qi Z, et al. Paint3D: Paint anything 3D with lighting-immersive texture generation. \textit{Proceedings of the IEEE/CVF Conference on Computer Vision and Pattern Recognition}, 2024.

\bibitem{cheng2024mvpaint}
Cheng W, Mu J, Zeng X, et al. MVPaint: Synchronized multi-view diffusion for painting anything 3D. \textit{Proceedings of the IEEE/CVF Conference on Computer Vision and Pattern Recognition}, 2025.

\bibitem{lai2025hunyuan3d}
Zhao Z, Lai Z, Lin Q, et al. Hunyuan3D 2.0: Scaling diffusion models for high resolution textured 3D assets generation. arXiv:2501.12202, 2025.

\bibitem{huang2024mvadapter}
Huang Z, et al. MV-Adapter: Multi-view consistent image generation made easy. \textit{Proceedings of the IEEE/CVF International Conference on Computer Vision}, 2025.

\bibitem{tang2023mvdiffusion}
Tang S, Zhang F, Chen J, Wang P, Furukawa Y. MVDiffusion: Enabling holistic multi-view image generation with correspondence-aware diffusion. \textit{Advances in Neural Information Processing Systems}, 2023.

\bibitem{liu2024synchronized}
Liu Y X, Xie M S, Liu H Y, Wong T T. Text-guided texturing by synchronized multi-view diffusion. \textit{SIGGRAPH Asia 2024 Conference Papers}, 2024.

\bibitem{xu2026wondertex}
Xu Q, Han X G, Zhang L. WonderTex: Consistent-and-seamless texture generation with text-guided multi-view image diffusion models. \textit{IEEE Transactions on Visualization and Computer Graphics}, 2026, 32(7): 5815--5825.

\bibitem{mou2023t2iadapter}
Mou C, Wang X, Xie L, et al. T2I-Adapter: Learning adapters to dig out more controllable ability for text-to-image diffusion models. \textit{Proceedings of the AAAI Conference on Artificial Intelligence}, 2024.

\bibitem{zhang2023controlnet}
Zhang L, Rao A, Agrawala M. Adding conditional control to text-to-image diffusion models. \textit{Proceedings of the IEEE/CVF International Conference on Computer Vision}, 2023.

\bibitem{ye2023ipadapter}
Ye H, Zhang J, Liu S, Han X, Yang W. IP-Adapter: Text compatible image prompt adapter for text-to-image diffusion models. arXiv:2308.06721, 2023.

\bibitem{hu2022lora}
Hu E J, Shen Y, Wallis P, et al. LoRA: Low-rank adaptation of large language models. \textit{International Conference on Learning Representations}, 2022.

\bibitem{liu2024dora}
Liu S Y, Wang C Y, Yin H, et al. DoRA: Weight-decomposed low-rank adaptation. \textit{International Conference on Machine Learning}, 2024.

\bibitem{zhang2023adalora}
Zhang Q, Zuo S, Liang C, et al. AdaLoRA: Adaptive budget allocation for parameter-efficient fine-tuning. \textit{International Conference on Learning Representations}, 2023.

\bibitem{kulkarni2026ssi}
Kulkarni A S, et al. Scheduled Style Injection: Expanding the Style-Content Pareto Frontier in Training-Free Diffusion-Based Style Transfer. \textit{Proceedings of the IEEE/CVF Conference on Computer Vision and Pattern Recognition Workshops (CVPRW), NTIRE Workshop}, 2026.

\bibitem{ho2020ddpm}
Ho J, Jain A, Abbeel P. Denoising diffusion probabilistic models. \textit{Advances in Neural Information Processing Systems}, 2020.

\bibitem{song2021ddim}
Song J, Meng C, Ermon S. Denoising diffusion implicit models. \textit{International Conference on Learning Representations}, 2021.

\bibitem{ho2022cfg}
Ho J, Salimans T. Classifier-free diffusion guidance. \textit{NeurIPS Workshop on Deep Generative Models and Downstream Applications}, 2021. arXiv:2207.12598.

\bibitem{castillo2023adaptive}
Castillo A, Kohler J, P\'erez J C, et al. Adaptive guidance: Training-free acceleration of conditional diffusion models. \textit{Proceedings of the AAAI Conference on Artificial Intelligence}, 2025.

\bibitem{wang2004ssim}
Wang Z, Bovik A C, Sheikh H R, Simoncelli E P. Image quality assessment: From error visibility to structural similarity. \textit{IEEE Transactions on Image Processing}, 2004, 13(4): 600--612.

\bibitem{zhang2018lpips}
Zhang R, Isola P, Efros A A, et al. The unreasonable effectiveness of deep features as a perceptual metric. \textit{CVPR}, 2018.

\end{thebibliography}
\endgroup
\end{document}
